We characterise what limits float32 statevector circuit simulation on \uno, \espA\ and \espB, and how well they can be predicted. For \hall, \bell\ and \qft\ the memory wall is at $n=\MemWallUno$, $\MemWallEspLo$ and $\MemWallEspHi$; the rule $\lfloor\log_2(B/\AmpBytes)\rfloor$ with $B$ read at $n=2$ is consistent with the largest working $n$ (on the \espB\ only with the largest free block, chosen after the walls were known), at a resolution of a factor of two. \grover\ on the ESP boards stopped at a time cap before any wall. A per-gate-kind time model fitted at small $n$ has a held-out error of \MapeTwoClean\,\% (MAPE, \NCleanTwo\ non-weak firmware-class-A points) against \MapeOneClean\,\% for one constant per board; \grover\ on the ESP boards is not scored, the model is not identifiable on the Uno at the primary split, was not defined blind, several rules were set after the first scores, and its constants are fits, not gate costs. Offloading to a PC overtakes local execution within brackets of measured $n$ in \CrossBoth\ of \NCross\ board--circuit pairs under both bounds; the brackets depend on the PC's simulation speed (up to \SimSpreadMax$\times$ spread between sessions). One unit per board, one PC, one baud rate.
